\begin{theorem}[Finite-chain normalization under nonzero proportionality]
    \label{thm:mpdo_finite_chain_normalized_proportionality}
    \lean{MPOTensor.normalizedMPO_eq_of_nonzeroProportionalMPV₂_at}
    \leanok
    \uses{def:mpo_tensor, def:mpo_to_mps, def:mpv}
    Let $M$ and $C$ be MPO tensors, possibly of different bond dimensions,
    and fix a chain length $N$. Write
    \begin{align}
        \widehat\rho^{(N)}(M)
        &=
        \tr\!\left[\rho^{(N)}(M)\right]^{-1}\rho^{(N)}(M),
        \notag
    \end{align}
    with the inverse of zero understood to be zero. Suppose that, for some
    $c\in\C\setminus\{0\}$, the doubled-index matrix-product vectors satisfy
    $\ket{V^{(N)}(M)}=c\ket{V^{(N)}(C)}$. Then
    \begin{align}
        \widehat\rho^{(N)}(M)
        &=
        \widehat\rho^{(N)}(C).
        \notag
    \end{align}

    The proportionality hypothesis is the finite-length form of the
    hypothesis in \cite[Theorem~2.10]{Cirac2016MPDO_arXiv}. The conclusion is
    an auxiliary cancellation used in the comparison preceding
    \cite[Appendix~C.2, lines~1597--1619]{Cirac2016MPDO_arXiv}.
\end{theorem}

\begin{proof}\leanok
    Reading a doubled physical index as a ket--bra pair gives
    $\rho^{(N)}(M)=c\rho^{(N)}(C)$, and hence
    $\tr[\rho^{(N)}(M)]=c\tr[\rho^{(N)}(C)]$. Since $c\ne0$,
    \begin{align}
        \bigl(c\tr[\rho^{(N)}(C)]\bigr)^{-1}
        c\rho^{(N)}(C)
        &=
        \tr[\rho^{(N)}(C)]^{-1}\rho^{(N)}(C).
        \notag
    \end{align}
    This identity also holds when the two proportional traces vanish.
\end{proof}

\begin{theorem}[Finite-chain reduced marginals under nonzero proportionality]
    \label{thm:mpdo_finite_chain_reduced_proportionality}
    \lean{MPOTensor.reducedBlockState_eq_of_nonzeroProportionalMPV₂_at}
    \leanok
    \uses{thm:mpdo_finite_chain_normalized_proportionality,
      def:mpo_to_mps}
    Under the hypotheses of
    Theorem~\ref{thm:mpdo_finite_chain_normalized_proportionality}, for every
    $L\leq N$ the reduced operators on the first $L$ sites agree:
    \begin{align}
        \tr_{N-L}\!\left[\widehat\rho^{(N)}(M)\right]
        &=
        \tr_{N-L}\!\left[\widehat\rho^{(N)}(C)\right].
        \notag
    \end{align}
\end{theorem}

\begin{proof}\leanok
    \uses{thm:mpdo_finite_chain_normalized_proportionality}
    Apply the same finite-index identification and partial trace to both
    sides of
    Theorem~\ref{thm:mpdo_finite_chain_normalized_proportionality}.
\end{proof}

\begin{theorem}[Eventual equality of normalized finite-chain operators]
    \label{thm:mpdo_eventual_normalized_proportionality}
    \lean{MPSTensor.EventuallyNonzeroProportionalMPV₂.%
      eventually_normalizedMPO_eq}
    \leanok
    \uses{def:eventually_nonzero_proportional_mpv, def:mpo_tensor,
      def:mpo_to_mps}
    If the doubled-index tensors of $M$ and $C$ are eventually nonzero
    proportional, then, for every sufficiently large $N$,
    \begin{align}
        \widehat\rho^{(N)}(M)
        &=
        \widehat\rho^{(N)}(C).
        \notag
    \end{align}
\end{theorem}

\begin{proof}\leanok
    \uses{thm:mpdo_finite_chain_normalized_proportionality}
    At each sufficiently large length, apply
    Theorem~\ref{thm:mpdo_finite_chain_normalized_proportionality} to the
    nonzero proportionality scalar at that length.
\end{proof}

\begin{theorem}[Eventual equality of fixed-size reduced marginals]
    \label{thm:mpdo_eventual_reduced_proportionality}
    \lean{MPSTensor.EventuallyNonzeroProportionalMPV₂.%
      eventually_reducedBlockState_eq}
    \leanok
    \uses{def:eventually_nonzero_proportional_mpv, def:mpo_to_mps}
    If the doubled-index tensors of $M$ and $C$ are eventually nonzero
    proportional, then for each fixed $L$ and every sufficiently large
    $N\geq L$,
    \begin{align}
        \tr_{N-L}\!\left[\widehat\rho^{(N)}(M)\right]
        &=
        \tr_{N-L}\!\left[\widehat\rho^{(N)}(C)\right].
        \notag
    \end{align}
    This equality alone does not identify the virtual sector families of
    $M$ and $C$, and therefore does not prove the rank-one separation or the
    Markov decomposition asserted in
    \cite[Appendix~C.2, lines~1597--1619]{Cirac2016MPDO_arXiv}.
\end{theorem}

\begin{proof}\leanok
    \uses{thm:mpdo_finite_chain_reduced_proportionality}
    At each sufficiently large length, apply
    Theorem~\ref{thm:mpdo_finite_chain_reduced_proportionality}.
\end{proof}

\begin{lemma}[Marginal stability under source zero correlation length]
    \label{lem:mpdo_source_zcl_marginal_stability}
    \lean{MPOTensor.%
      reducedBlockState_apply_eq_trace_evalWord_mul_verticalLoop_pow}
    \lean{MPOTensor.%
      reducedBlockState_succ_succ_eq_succ_of_isSourceZCL}
    \leanok
    \uses{def:mpo_source_zcl, lem:mpdo_trace_loop_chain}
    Let $M$ have source zero correlation length. For every $N\geq0$, the
    normalized $N$-site marginal obtained by tracing two sites from
    $\sigma^{(N+2)}(M)$ equals the normalized $N$-site marginal obtained by
    tracing one site from $\sigma^{(N+1)}(M)$.

    This is the first replacement at
    \cite[Appendix~C.2, line~1613]{Cirac2016MPDO_arXiv}. It does not imply
    the rank-one trace factorization asserted later on that line.
\end{lemma}

\begin{proof}\leanok
    Write $E=\mathcal T_M$ and choose $\lambda>0$ such that
    $E^2=\lambda E$. For physical configurations $u,v$ of length $N$, the
    two marginal matrix elements are
    $\tr(M^{u,v}E^2)/\tr(E^{N+2})$ and
    $\tr(M^{u,v}E)/\tr(E^{N+1})$.
    The numerator and denominator on the left are respectively $\lambda$
    times those on the right, and source zero correlation length makes the
    denominator on the right nonzero.
\end{proof}

% **Local fix (additional marginal replacement):** the printed proof returns
% to a one-step coefficient after its first replacement.  See
% docs/paper-gaps/cpsv16_commuting_form_to_sal.tex.
\begin{theorem}[Iterated marginal stability under zero correlation length]
    \label{thm:mpdo_source_zcl_iterated_marginal_stability}
    \lean{MPOTensor.reducedBlockState_add_three_eq_succ_of_isSourceZCL}
    \leanok
    \uses{lem:mpdo_source_zcl_marginal_stability}
    Let $M$ have source zero correlation length. For every $L\geq0$, the
    normalized $L$-site marginal obtained by tracing three sites from
    $\sigma^{(L+3)}(M)$ equals the normalized $L$-site marginal obtained by
    tracing one site from $\sigma^{(L+1)}(M)$.
\end{theorem}

\begin{proof}\leanok
    First retain $L+1$ sites and apply
    Lemma~\ref{lem:mpdo_source_zcl_marginal_stability}. Trace the last of
    these retained sites, and then apply the same lemma while retaining $L$
    sites.
\end{proof}

% Source: arXiv:1606.00608, Appendix C.2, Proposition `4to2`, lines
% 1606--1617.
% **Local fix (adjacent marginal comparison):** the arbitrary positive suffix
% length includes the one-, two-, and three-site contractions used below. See
% docs/paper-gaps/cpsv16_commuting_form_to_sal.tex.
\begin{definition}[Fixed-sector suffix contraction]
    \label{def:mpdo_suffix_sector_contraction}
    \lean{MPOTensor.PhysicalSectorFactorization.suffixSectorContraction}
    \lean{MPOTensor.PhysicalSectorFactorization.%
      oneSuffixSectorContraction}
    \lean{MPOTensor.PhysicalSectorFactorization.%
      twoSuffixSectorContraction}
    \lean{MPOTensor.PhysicalSectorFactorization.%
      threeSuffixSectorContraction}
    \leanok
    \uses{def:mpdo_physical_sector_chain_coordinates}
    Let $S$ be the finite set of physical sectors and let
    $k=(k_0,\ldots,k_{L-1})\in S^L$ be a retained sector word. For
    $R\geq1$ and $x,y\in I_k$, define
    \begin{align}
        \Gamma_k^{(R)}(x,y)
        &=
        \sum_{t\in S^R}\ \sum_{z\in I_t}
        \Omega_{kt}(xz,yz),
        \notag
    \end{align}
    where $kt$ and $xz$ denote concatenation of sector words and their
    corresponding fiber coordinates. The cases $R=1,2,3$ are denoted by
    $\Gamma_k^{(1)}$, $\Gamma_k^{(2)}$, and $\Gamma_k^{(3)}$.
\end{definition}

\begin{theorem}[Contraction of a sector-chain fiber]
    \label{thm:mpdo_sector_chain_fiber_contraction}
    \lean{MPOTensor.PhysicalSectorFactorization.%
      sum_sectorChainFiber_neighboringOperator_diag}
    \leanok
    \uses{def:mpdo_physical_sector_chain_coordinates}
    Let $t=(t_0,\ldots,t_m)$ be a nonempty sector word. For functions
    $a$ on the left coordinate of $t_0$ and $b$ on the right coordinate of
    $t_m$,
    \begin{align}
      &\sum_{z\in I_t}
        a(z_0^{\mathrm L})
        \prod_{i=0}^{m-1}
          \eta_{t_i,t_{i+1}}
          \bigl((z_i^{\mathrm R},z_{i+1}^{\mathrm L}),
                (z_i^{\mathrm R},z_{i+1}^{\mathrm L})\bigr)
        b(z_m^{\mathrm R})
        \notag\\
      &\qquad=
        \left(\sum_u a(u)\right)
        \left(\prod_{i=0}^{m-1}\tr(\eta_{t_i,t_{i+1}})\right)
        \left(\sum_v b(v)\right).
    \end{align}
\end{theorem}

\begin{proof}\leanok
    Write $z_i=(\ell_i,r_i)$. Finite distributivity separates the left and
    right coordinates:
    \begin{align}
      &\sum_{\ell,r}
        a(\ell_0)
        \prod_{i=0}^{m-1}
          \eta_{t_i,t_{i+1}}
          \bigl((r_i,\ell_{i+1}),(r_i,\ell_{i+1})\bigr)
        b(r_m)
        \notag\\
      &\qquad=
        \left(\sum_{\ell_0}a(\ell_0)\right)
        \prod_{i=0}^{m-1}
          \left(\sum_{r_i,\ell_{i+1}}
            \eta_{t_i,t_{i+1}}
            \bigl((r_i,\ell_{i+1}),(r_i,\ell_{i+1})\bigr)\right)
        \left(\sum_{r_m}b(r_m)\right).
    \end{align}
    The internal double sum is
    $\tr(\eta_{t_i,t_{i+1}})$ for each $i$, which gives the asserted
    factorization.
\end{proof}

\begin{theorem}[Contraction of a nonempty suffix fiber]
    \label{thm:mpdo_nonempty_suffix_fiber_contraction}
    \lean{MPOTensor.PhysicalSectorFactorization.%
      sum_suffixFiber_cyclicNeighboringProduct}
    \leanok
    \uses{thm:mpdo_sector_chain_fiber_contraction}
    Let $k=(k_0,\ldots,k_n)$ be a nonempty retained sector word and
    $t=(t_0,\ldots,t_m)$ a nonempty suffix sector word. In open-edge
    coordinates $x=(x_{\mathrm{int}},x_{\partial})$ and
    $y=(y_{\mathrm{int}},y_{\partial})$,
    \begin{align}
      \sum_{z\in I_t}\Omega_{kt}(xz,yz)
      &=
      \left(
        \prod_{j=0}^{n-1}\eta_{k_j,k_{j+1}}
      \right)(x_{\mathrm{int}},y_{\mathrm{int}})
      \left(
        \tr_{\mathrm R}(\eta_{k_n,t_0})
        \otimes
        \tr_{\mathrm L}(\eta_{t_m,k_0})
      \right)(x_{\partial},y_{\partial})
      \prod_{i=0}^{m-1}\tr(\eta_{t_i,t_{i+1}}).
    \end{align}
\end{theorem}

\begin{proof}\leanok
    Split the cyclic product into the retained internal edges, the two
    boundary edges, and the internal suffix edges. Apply
    Theorem~\ref{thm:mpdo_sector_chain_fiber_contraction} to the suffix
    coordinates. The two boundary sums are the indicated partial traces.
\end{proof}

% Source: arXiv:1606.00608, Appendix C.2, Proposition `4to2`, lines
% 1606--1617.
% **Local fix (adjacent marginal comparison):** the cases R=1 and R=2 give
% the adjacent marginals whose boundary traces have coefficients T_C^2 and
% T_C^3. See docs/paper-gaps/cpsv16_commuting_form_to_sal.tex.
\begin{theorem}[Sector decomposition of a suffix marginal]
    \label{thm:mpdo_suffix_marginal_block_expansion}
    \lean{MPOTensor.PhysicalSectorFactorization.%
      reindex_reducedBlockState_add_eq_suffixSectorContraction}
    \lean{MPOTensor.PhysicalSectorFactorization.%
      reindex_reducedBlockState_add_one_eq_oneSuffixSectorContraction}
    \lean{MPOTensor.PhysicalSectorFactorization.%
      reindex_reducedBlockState_add_two_eq_twoSuffixSectorContraction}
    \lean{MPOTensor.PhysicalSectorFactorization.%
      reindex_reducedBlockState_add_three_eq_threeSuffixSectorContraction}
    \leanok
    \uses{def:mpdo_suffix_sector_contraction,
      def:mpdo_physical_sector_chain_coordinates}
    Let $\widetilde{\mathcal K}$ be the tensor in physical-sector
    coordinates. For $L\geq0$ and $R\geq1$, the normalized $L$-site
    marginal obtained from a chain of length $L+R$ satisfies
    \begin{align}
        \reindex_{\Phi_L}\!
        \left(
          \tr_R\!\left[\widehat\rho^{(L+R)}
          (\widetilde{\mathcal K})\right]
        \right)
        &=
        \tr\!\left[\rho^{(L+R)}
          (\widetilde{\mathcal K})\right]^{-1}
        \bigoplus_{k\in S^L}\Gamma_k^{(R)}.
        \notag
    \end{align}
    In particular, this identity holds for $R=1$, $R=2$, and $R=3$.
\end{theorem}

\begin{proof}\leanok
    \uses{thm:mpdo_physical_sector_chain_decomposition}
    Decompose each discarded physical configuration uniquely as a suffix
    sector word $t\in S^R$ and a fiber coordinate $z\in I_t$. By
    Theorem~\ref{thm:mpdo_physical_sector_chain_decomposition}, matrix
    elements with distinct retained sector words vanish. For a fixed
    retained word $k$, summing the diagonal entries of the discarded fiber
    gives
    \begin{align}
        \sum_{t\in S^R}\ \sum_{z\in I_t}
        \Omega_{kt}(xz,yz)
        &=\Gamma_k^{(R)}(x,y).
        \notag
    \end{align}
    Multiplication by the inverse trace of the length-$L+R$ operator gives
    the stated normalization.
\end{proof}

\begin{definition}[Rank-one trace factorization for neighboring operators]
    \label{def:mpdo_generic_eta_rank_one_trace_factorization}
    \lean{Matrix.EtaRankOneTraceFactorization}
    \leanok
    A rank-one trace factorization of neighboring matrices $\eta_{q,h}$
    consists of real numbers $a_q,b_h$ such that
    \begin{align}
        \tr(\eta_{q,h})&=a_qb_h,
        \notag\\
        \sum_q a_qb_q&=1.
        \notag
    \end{align}
    The factorization is an explicit assumption; it is not deduced from zero
    correlation length. It is the factorization invoked at
    \cite[Appendix~C.2, line~1613]{Cirac2016MPDO_arXiv}.
    % Scope restriction documented in
    % docs/paper-gaps/cpsv16_commuting_form_to_sal.tex.
\end{definition}

\begin{lemma}[Separation of the two-boundary edge contraction]
    \label{lem:mpdo_eta_two_boundary_contraction_separated}
    \lean{Matrix.etaTwoBoundaryContraction_eq_separated}
    \leanok
    \uses{def:mpdo_generic_eta_rank_one_trace_factorization,
        def:mpdo_right_partial_trace_map, def:partial_trace_left}
    Suppose $\eta_{q,h}$ has a rank-one trace factorization. For fixed
    boundary sectors $u,v$, set
    \begin{align}
        R_{u,q}&=\tr_{H_{q,l}}(\eta_{u,q}),
        \notag\\
        L_{h,v}&=\tr_{H_{h,r}}(\eta_{h,v}).
        \notag
    \end{align}
    Then
    \begin{align}
        \sum_{q,h}\tr(\eta_{q,h})(R_{u,q}\otimes L_{h,v})
        &=
        \left(\sum_q a_qR_{u,q}\right)\otimes
        \left(\sum_h b_hL_{h,v}\right).
        \notag
    \end{align}
    This is the corrected edge contraction at
    \cite[Appendix~C.2, lines~1613--1617]{Cirac2016MPDO_arXiv}.
\end{lemma}

\begin{proof}\leanok
    \uses{def:mpdo_generic_eta_rank_one_trace_factorization}
    Substitute $\tr(\eta_{q,h})=a_qb_h$ and distribute the two finite sums.
\end{proof}

% **Scope restriction (rank-one trace factorization):** the theorem below
% assumes tr(eta_{q,h}) = a_q b_h.  It does not derive that identity from zero
% correlation length and does not prove Proposition 4to2, a Markov
% decomposition, or saturation of the area law.  Source line 1616 suppresses
% the partial traces on the two surviving boundary operators.  Documented in
% docs/paper-gaps/cpsv16_commuting_form_to_sal.tex.
\begin{theorem}[Conditional fourth-region boundary trace]
    \label{thm:mpdo_eta_fourth_region_trace_formula}
    \lean{Matrix.eta_fourth_region_trace_formula}
    \leanok
    \uses{def:mpdo_generic_eta_rank_one_trace_factorization,
        def:mpdo_right_partial_trace_map, def:partial_trace_left}
    Let $\eta_{q,h}$ be arbitrary matrices on
    $H_{q,r}\otimes H_{h,l}$ with a rank-one trace factorization. For a
    nonempty retained sector word $k=(k_0,\ldots,k_{L-1})$, write
    \begin{align}
        P_k&=\bigotimes_{n=0}^{L-2}\eta_{k_n,k_{n+1}},
        \notag\\
        R_{u,q}&=\tr_{H_{q,l}}(\eta_{u,q}),
        \notag\\
        L_{h,v}&=\tr_{H_{h,r}}(\eta_{h,v}).
        \notag
    \end{align}
    The edge contraction corresponding to tracing the two boundary sites of
    the longer cyclic chain is
    \begin{align}
        &\bigoplus_k P_k\otimes
          \sum_{q,h}\tr(\eta_{q,h})
            (R_{k_{L-1},q}\otimes L_{h,k_0})
        \notag\\
        &\quad=
          \bigoplus_k P_k\otimes
          \left(\sum_q a_qR_{k_{L-1},q}\right)\otimes
          \left(\sum_h b_hL_{h,k_0}\right).
        \notag
    \end{align}
    This is the edge-coordinate interpretation of the displayed
    fourth-region trace in
    \cite[Appendix~C.2, lines~1606 and 1613--1617]{Cirac2016MPDO_arXiv}.
\end{theorem}

\begin{proof}\leanok
    \uses{lem:mpdo_eta_two_boundary_contraction_separated}
    Apply Lemma~\ref{lem:mpdo_eta_two_boundary_contraction_separated} for
    each retained sector word. Tensor with the unchanged product $P_k$ and
    take the direct sum over $k$.
\end{proof}

\begin{remark}
    The theorem applies, in particular, to the positive chain-independent
    neighboring family supplied by
    Theorem~\ref{thm:mpdo_eta_local_sigma_nk2}, once the rank-one trace
    factorization is given. It does not derive that factorization, identify
    this two-boundary contraction with a one-site marginal, or prove a Markov
    decomposition or saturation of the area law.
\end{remark}

\begin{theorem}[A nonminimal sector factorization need not be primitive]
    \label{thm:mpdo_commuting_bond_trace_matrix_primitivity_obstruction}
    \lean{MPOTensor.CommutingBondTraceMatrixObstruction.%
      normal_sourceZCL_fixed_commutingBond_does_not_force_traceMatrix_primitive}
    \leanok
    \uses{def:mpdo_eta_local_structure_data, def:mpo_source_zcl, def:mpdo,
      def:injective, def:normal}
    There is an injective normal tensor $\mathcal K$, of physical dimension two
    and bond dimension one, which generates MPDOs and has source zero
    correlation length, together with one positive two-site bond $B$ whose
    translates commute. For every $N\geq2$,
    \begin{align}
        \sigma^{(N)}(\mathcal K)
        &=\prod_{i=0}^{N-1}B_{i,i+1},
        \notag
    \end{align}
    while the trace matrix of a chosen two-sector Beigi factorization is
    \begin{align}
        T&=\begin{pmatrix}1&0\\0&0\end{pmatrix}.
        \notag
    \end{align}
    In particular, $T$ is not primitive.

    This shows that primitivity cannot be inferred for an arbitrary selected
    factorization in the step suggested at
    \cite[Appendix~C.2, line~1613]{Cirac2016MPDO_arXiv}. The same tensor also
    has a one-sector factorization with positive semidefinite neighboring
    operator and primitive trace matrix. Hence the example does not exclude a
    source-faithful minimal or visible-sector selection, and it does not
    disprove the stated commuting-form-to-SAL implication.
\end{theorem}

\begin{proof}\leanok
    Let $\mathcal K^{0,0}=1$ and let all other physical entries vanish. The
    matrices $\{\mathcal K^{i,j}\}_{i,j=0}^1$ span $M_1(\mathbb C)$. For every
    $X\in M_1(\mathbb C)$, the doubled-index transfer $\mathcal E_{\mathcal K}$
    and the physical-trace transfer $P_{\mathcal K}$ satisfy
    \begin{align}
        \mathcal E_{\mathcal K}(X)&=X,
        \notag\\
        P_{\mathcal K}&=I_1,
        \notag\\
        P_{\mathcal K}^2&=P_{\mathcal K}.
        \notag
    \end{align}
    Hence $\mathcal K$ is injective and normal and has source zero correlation
    length. Decompose the physical space nonminimally as
    $\mathbb C\oplus\mathbb C$. With $w_0=1$ and $w_1=0$, take both sector
    factors to be the scalar $w_k$ in sector $k$. The neighboring operators
    are
    \begin{align}
        \eta_{k,h}
        &=w_kw_h I_1
          =\delta_{k,0}\delta_{h,0}I_1
          \succeq0.
        \notag
    \end{align}
    In the sector coordinates the corresponding physical bond is
    \begin{align}
        B
        &=\sum_{k,h}\Id_{H_L^{(k)}}\otimes\eta_{k,h}\otimes
          \Id_{H_R^{(h)}}
        \notag\\
        &=\Id_{H_L^{(0)}}\otimes I_1\otimes\Id_{H_R^{(0)}}.
        \notag
    \end{align}
    Since only the all-zero physical-index word contributes, the
    factorization identity is
    \begin{align}
        \sigma^{(N)}(\mathcal K)
        &=\prod_{i=0}^{N-1}B_{i,i+1}.
        \notag
    \end{align}
    Consequently,
    \begin{align}
        T_{k,h}
        &=\tr(\eta_{k,h})
          =\delta_{k,0}\delta_{h,0},
        \notag\\
        T&=\begin{pmatrix}1&0\\0&0\end{pmatrix}.
        \notag
    \end{align}
    For every $n\geq1$,
    $T^n=\begin{pmatrix}1&0\\0&0\end{pmatrix}$.
    Thus no positive power has all entries strictly positive, so $T$ is not
    primitive. For the alternative factorization, take one sector
    $B^L=\mathbb C^2$ and $B^R=\mathbb C$, with
    \begin{align}
        l&=\begin{pmatrix}1&0\\0&0\end{pmatrix},
        \notag\\
        r&=(1).
        \notag
    \end{align}
    Its sole neighboring operator is
    $\widehat\eta_{0,0}=\diag(1,0)\succeq0$, and hence its trace
    matrix is $\widehat T=(\tr(\widehat\eta_{0,0}))=(1)$.
    Every positive power of $\widehat T$ equals $(1)$, so this one-sector
    trace matrix is primitive.
\end{proof}

% **Local fix (cyclic-active restriction):** CPSV16 Lemma `propSN` assumes SAL.
% The following restricted construction repairs the arbitrary-factorization obstruction.
% See docs/paper-gaps/cpsv16_commuting_form_to_sal.tex.
\begin{definition}[Cyclic-active sectors]
    \label{def:mpdo_cyclic_active_sector}
    \lean{MPOTensor.PhysicalSectorFactorization.IsCyclicActiveSector}
    \lean{MPOTensor.PhysicalSectorFactorization.cyclicActiveWeight}
    \lean{MPOTensor.PhysicalSectorFactorization.cyclicActiveWeight_ne_zero_iff}
    \lean{MPOTensor.PhysicalSectorFactorization.CyclicActiveSector}
    \lean{MPOTensor.PhysicalSectorFactorization.cyclicActiveSectorTraceMatrix}
    \lean{MPOTensor.PhysicalSectorFactorization.neighboringOperator_trace_re_pos_iff}
    \leanok
    \uses{def:mpdo_physical_sector_factorization,
      def:mpdo_minimal_neighboring_operator}
    For a physical-sector factorization with neighboring operators
    $\eta_{k,h}$, retain a sector $k$ when there is a sector $h$ such that
    $\eta_{k,h}\ne0$ and $h\leadsto k$, where $\leadsto$ denotes reachability
    by nonzero neighboring operators.
    The first edge $k\to h$ is explicit, so the resulting closed walk has
    positive length even when the return path is empty. The set $C$ of
    retained sectors is therefore exactly the set of vertices occurring in
    positive-length directed cycles, in the sense of
    \cite[Section~III]{Beigi2012Classification}. Define, for $k,h\in C$,
    \begin{align}
        T_{k,h}&=\re\tr(\eta_{k,h}).
        \notag
    \end{align}
    For positive semidefinite neighboring operators the trace is real, so
    $T_{k,h}=\tr(\eta_{k,h})$, and $T_{k,h}>0$ is equivalent to
    $\eta_{k,h}\ne0$.
\end{definition}

% **Local fix (cyclic-active restriction):** this theorem concerns the restricted
% cyclic support, not the full trace matrix in CPSV16 Lemma `propSN`. See
% docs/paper-gaps/cpsv16_commuting_form_to_sal.tex.
\begin{theorem}[Deletion of sectors absent from cyclic products]
    \label{thm:mpdo_cyclic_active_deletion}
    \lean{MPOTensor.PhysicalSectorFactorization.isCyclicActiveSector_of_cyclic_edges}
    \lean{MPOTensor.PhysicalSectorFactorization.neighboringOperator_ne_zero_of_cyclicNeighboringProduct_ne_zero}
    \lean{MPOTensor.PhysicalSectorFactorization.isCyclicActiveSector_of_cyclicNeighboringProduct_ne_zero}
    \lean{MPOTensor.PhysicalSectorFactorization.cyclicNeighboringProduct_eq_zero_of_not_isCyclicActiveSector}
    \lean{MPOTensor.PhysicalSectorFactorization.cyclicNeighboringProduct_eq_zero_of_not_forall_isCyclicActiveSector}
    \lean{MPOTensor.PhysicalSectorFactorization.%
      map_cyclicNeighboringProduct_eq_zero_of_not_forall_isCyclicActiveSector}
    \leanok
    \uses{def:mpdo_cyclic_active_sector,
      def:mpdo_physical_sector_chain_coordinates}
    Let $N\geq1$ and let $k=(k_0,\ldots,k_{N-1})$ be a cyclic sector
    configuration. If $k_n\notin C$ for some $n$, then $\Omega_k=0$.
    Thus deleting the sectors outside $C$ leaves every finite cyclic
    neighboring product unchanged. For every linear contraction $\Phi$,
    including a composition of reindexings and partial traces,
    $\Omega_k=0$ implies $\Phi(\Omega_k)=0$.
\end{theorem}

\begin{proof}\leanok
    If $\Omega_k\ne0$, each factor $\eta_{k_n,k_{n+1}}$ is nonzero; otherwise
    every matrix entry of the product vanishes. Starting with the edge
    $k_n\to k_{n+1}$ and following the remaining cyclic edges gives a return
    path to $k_n$. Hence every $k_n$ belongs to $C$. The last assertion is
    the linear identity $\Phi(0)=0$.
\end{proof}

% Source: arXiv:1606.00608, Appendix C.2, Proposition `4to2`, lines
% 1597--1619.
% **Scope restriction (selected-sector visibility):** the argument remains in the
% auxiliary space of the original injective tensor and transfers only literal
% zero edges from the selected factorization. See
% docs/paper-gaps/cpsv16_commuting_form_to_sal.tex.
\begin{theorem}[Trace-visible recurrence of selected fixed-product sectors]
    \label{thm:mpdo_selected_fixed_product_cyclic_recurrence}
    \lean{MPOTensor.EtaLocalStructureData.%
      exists_positive_scalar_mpo_changePhysicalBasis_eq_smul_selected}
    \lean{MPOTensor.PhysicalSectorFactorization.%
      changePhysicalBasis_physicalCoordinateMatrix_isInjective}
    \lean{MPOTensor.PhysicalSectorFactorization.%
      trace_sectorCoordinateTensor_mul_eq_zero_of_not_isCyclicActiveSector}
    \lean{MPOTensor.PhysicalSectorFactorization.%
      changePhysicalBasis_apply_ne_eq_zero_of_mpo_two_eq_smul}
    \lean{MPOTensor.PhysicalSectorFactorization.%
      changePhysicalBasis_apply_eq_zero_of_not_isCyclicActiveSector}
    \lean{MPOTensor.PhysicalSectorFactorization.%
      changePhysicalBasis_mul_eq_zero_of_neighboringOperator_eq_zero}
    \lean{MPOTensor.PhysicalSectorFactorization.retainedBulkProduct}
    \lean{MPOTensor.PhysicalSectorFactorization.%
      exists_retainedBulkProduct_diag_ne_zero_of_neighboringOperator_ne_zero}
    \lean{MPOTensor.PhysicalSectorFactorization.%
      exists_cyclicNeighboringProduct_diag_ne_zero_of_neighboringOperator_ne_zero}
    \lean{MPOTensor.PhysicalSectorFactorization.%
      exists_retainedBulkProduct_diag_ne_zero_of_isCyclicActiveSector}
    \lean{MPOTensor.PhysicalSectorFactorization.%
      exists_cyclicNeighboringProduct_diag_ne_zero_of_isCyclicActiveSector}
    \lean{MPOTensor.PhysicalSectorFactorization.%
      exists_changePhysicalBasis_apply_ne_zero_of_isCyclicActiveSector}
    \lean{MPOTensor.PhysicalSectorFactorization.%
      CyclicActiveOriginalEntryIndex}
    \lean{MPOTensor.PhysicalSectorFactorization.%
      cyclicActiveOriginalOneSiteMatrixFamily}
    \lean{MPOTensor.PhysicalSectorFactorization.%
      cyclicActiveOriginalOneSiteMatrixFamily_span_eq_top}
    \lean{MPOTensor.PhysicalSectorFactorization.%
      cyclicActive_reflTransGen_neighboringOperator_ne_zero_of_closed_mpo}
    \lean{MPOTensor.EtaLocalStructureData.%
      selectedFixedProduct_cyclicActive_reflTransGen_neighboringOperator_ne_zero}
    \leanok
    \uses{def:mpdo, def:injective, def:normal, def:mpo_source_zcl,
      def:mpdo_eta_local_structure_data,
      thm:mpdo_eta_fixed_bond_product_tensor,
      thm:mpdo_fixed_bond_positive_physical_sector_representative,
      def:mpdo_physical_sector_factorization,
      def:mpdo_cyclic_active_sector}
    Let $K$ be an injective normal matrix-product density tensor with source
    zero correlation length and an eta-local structure. Let $C$ be the exact
    fixed-product tensor selected from its commuting bond, and let $F$ be the
    associated positive physical-sector factorization. Then any two sectors
    in the positive-length cyclic support of $F$ are joined by a directed path
    of nonzero neighboring operators.
\end{theorem}

\begin{proof}\leanok
    At lengths two and three, express the closed operators of $K$ in the
    physical coordinates of $F$. They are nonzero scalar multiples of the
    corresponding closed operators of $C$. The length-two identity and
    nondegeneracy of the trace pairing show that cross-sector matrices of
    $K$ vanish, as do diagonal matrices belonging to sectors outside cyclic
    support. Hence the remaining cyclic-sector matrices span the full
    auxiliary matrix algebra.

    Positivity supplies a nonzero diagonal coefficient on a closed sector
    word through every cyclic sector. Its proportional coefficient for $K$
    shows that each cyclic sector contributes a nonzero one-site matrix of
    $K$. The length-three identity transfers a zero neighboring operator of
    $F$ to a zero ordered product of the corresponding two sector spaces of
    $K$.

    Fix a cyclic sector and let $S$ be the sectors reachable from it. If a
    cyclic sector lay outside $S$, the matrix spaces generated by $S$ and by
    its complement would both be nonzero, their sum would be the full matrix
    algebra, and every product from the first space to the second would
    vanish. This contradicts the one-sided matrix-space obstruction.
\end{proof}

\begin{theorem}[Selected fixed-product reduced states]
    \label{thm:mpdo_selected_fixed_product_reduced_states}
    \lean{MPOTensor.EtaLocalStructureData.%
      selectedFixedProduct_reducedBlockState_eq_changePhysicalBasis}
    \leanok
    \uses{def:mpdo_eta_local_structure_data,
      thm:mpdo_fixed_bond_positive_physical_sector_representative}
    Let $K$ have an eta-local structure, let $F$ be the positive
    physical-sector factorization of its selected fixed-product tensor, and
    express $K$ in the physical coordinates of $F$. For every $N\geq2$ and
    $L\leq N$, the selected sector-coordinate tensor and the transformed
    source tensor have the same normalized $L$-site reduced state.
\end{theorem}

\begin{proof}\leanok
    \uses{thm:mpdo_finite_chain_normalized_proportionality}
    At length $N$, the unnormalized periodic operators satisfy
    \[
      \rho^{(N)}(K_F)=c_N\,\rho^{(N)}(C_F),
      \qquad c_N>0.
    \]
    The scalar cancels under trace normalization, and taking the prescribed
    partial trace preserves the equality.
\end{proof}

\begin{theorem}[Adjacent selected fixed-product marginals]
    \label{thm:mpdo_selected_fixed_product_adjacent_marginals}
    \lean{MPOTensor.EtaLocalStructureData.%
      selectedFixedProduct_reducedBlockState_succ_succ_eq_succ_of_isSourceZCL}
    \leanok
    \uses{def:mpo_source_zcl, def:mpdo_eta_local_structure_data,
      thm:mpdo_selected_fixed_product_reduced_states}
    Let $K$ have an eta-local structure and source zero correlation length,
    and let $C$ be the selected fixed-product tensor in its physical-sector
    coordinates. For every $L>0$, the normalized $L$-site reduced states of
    $C$ obtained from chain lengths $L+2$ and $L+1$ are equal.
\end{theorem}

\begin{proof}\leanok
    In the selected physical coordinates, the periodic operators of $K$ at
    both lengths are positive nonzero scalar multiples of those of $C$.
    Cancel these scalars by
    Theorem~\ref{thm:mpdo_selected_fixed_product_reduced_states}, transport
    through the physical isometry, and apply source marginal stability.
\end{proof}

\begin{theorem}[Iterated selected fixed-product marginals]
    \label{thm:mpdo_selected_fixed_product_iterated_marginals}
    \lean{MPOTensor.EtaLocalStructureData.%
      selectedFixedProduct_reducedBlockState_add_three_eq_succ_of_isSourceZCL}
    \leanok
    \uses{def:mpo_source_zcl, def:mpdo_eta_local_structure_data,
      thm:mpdo_selected_fixed_product_reduced_states,
      thm:mpdo_source_zcl_iterated_marginal_stability}
    Under the preceding hypotheses, for every $L>0$, the normalized
    $L$-site reduced states of the selected fixed-product tensor obtained
    from chain lengths $L+3$ and $L+1$ are equal.
\end{theorem}

\begin{proof}\leanok
    Apply the two-step source marginal identity to $K$, transport both sides
    through the selected physical isometry, and cancel the two positive
    realization scalars.
\end{proof}

\begin{theorem}[Positive traces of the selected fixed product]
    \label{thm:mpdo_selected_fixed_product_positive_traces}
    \lean{MPOTensor.EtaLocalStructureData.%
      exists_pos_trace_mpo_selectedFixedProduct_of_isSourceZCL}
    \leanok
    \uses{def:mpo_source_zcl, def:mpdo_eta_local_structure_data,
      thm:mpdo_fixed_bond_positive_physical_sector_representative}
    Let $K$ have an eta-local structure and source zero correlation length.
    For every $N\geq2$, the trace of the selected fixed-product periodic
    operator is a strictly positive real number.
\end{theorem}

\begin{proof}\leanok
    The selected neighboring operators are positive semidefinite, so their
    cyclic product is positive semidefinite. Its trace cannot vanish:
    positive proportionality would then force the trace of the source
    periodic operator to vanish, contrary to source zero correlation length.
    A nonzero positive semidefinite matrix has strictly positive real trace.
\end{proof}

\begin{theorem}[Stabilization of the support of an irreducible square]
    \label{thm:matrix_irreducible_square_support_stabilization}
    \lean{Matrix.IsIrreducible.%
      pow_two_pos_of_pow_two_support_eq_pow_three_support}
    \leanok
    Let $T$ be a nonnegative irreducible real matrix on a finite index set.
    Suppose that
    \begin{align}
        (T^2)_{i,j}>0
        \quad\Longleftrightarrow\quad
        (T^3)_{i,j}>0
        \notag
    \end{align}
    for every pair $i,j$. Then $(T^2)_{i,j}>0$ for every pair $i,j$.
\end{theorem}

\begin{proof}\leanok
    Right multiplication by $T$ propagates equality of the positive supports
    of two consecutive powers. Hence all powers from the second onward have
    the same positive support. Irreducibility supplies a positive path from
    $i$ to $j$. If its length is one, append a positive return path at $j$.
    The resulting path has length at least two, so support stabilization
    reduces it to a positive path of length two from $i$ to $j$.
\end{proof}

% Source: arXiv:1606.00608, Appendix C.2, Proposition `4to2`, lines
% 1606--1617.
% **Scope restriction (coefficient extraction):** recurrence, adjacent
% marginal equality, and positive finite-chain normalizations are explicit
% hypotheses. See docs/paper-gaps/cpsv16_commuting_form_to_sal.tex.
\begin{theorem}[Strict positivity from adjacent cyclic-sector marginals]
    \label{thm:mpdo_cyclic_active_square_positive_adjacent_marginals}
    \lean{MPOTensor.PhysicalSectorFactorization.%
      cyclicActiveSectorTraceMatrix_pow_two_pos_of_adjacent_marginals}
    \leanok
    \uses{def:mpdo_physical_sector_factorization,
      def:mpdo_cyclic_active_sector}
    Let $F$ be a physical-sector factorization whose neighboring operators are
    positive semidefinite, and let $T_C$ be the trace matrix on its
    cyclic-active sectors. Assume that every ordered pair of cyclic-active
    sectors is joined by a path of nonzero neighboring operators. Assume also
    that, for every $L>0$, the normalized $L$-site marginals obtained from
    chain lengths $L+2$ and $L+1$ are equal, and that the closed operator has
    a positive real trace at every length $N\geq2$. Then
    $(T_C^2)_{q,h}>0$ for every pair of cyclic-active sectors $q,h$.
\end{theorem}

\begin{proof}\leanok
    \uses{def:mpdo_suffix_sector_contraction,
      thm:mpdo_suffix_marginal_block_expansion,
      thm:matrix_irreducible_square_support_stabilization}
    Fix $q,h$ and choose a retained sector path from $h$ to $q$. Positivity
    makes its retained trace product $P_{h,q}$ nonzero. Tracing one suffix
    site and two suffix sites gives the respective block-trace coefficients
    $P_{h,q}(T_C^2)_{q,h}$ and $P_{h,q}(T_C^3)_{q,h}$. Equality of the
    normalized marginals and positivity of their normalizations show that
    $(T_C^2)_{q,h}$ vanishes exactly when $(T_C^3)_{q,h}$ vanishes.
    Recurrence makes $T_C$ irreducible, so
    Theorem~\ref{thm:matrix_irreducible_square_support_stabilization} applies.
\end{proof}

% Source: arXiv:1606.00608, Appendix C.2, Proposition `4to2`, lines
% 1606--1617.
% **Scope restriction (coefficient extraction):** recurrence, adjacent
% marginal equality, and positive finite-chain normalizations are explicit
% hypotheses. See docs/paper-gaps/cpsv16_commuting_form_to_sal.tex.
\begin{theorem}[Normalization from adjacent cyclic-sector marginals]
    \label{thm:mpdo_cyclic_active_square_cube_adjacent_marginals}
    \lean{MPOTensor.PhysicalSectorFactorization.%
      exists_normalized_cyclicActiveSectorTraceMatrix_pow_two_eq_pow_three_of_adjacent_marginals}
    \leanok
    \uses{def:mpdo_physical_sector_factorization,
      def:mpdo_cyclic_active_sector}
    Let $F$ be a physical-sector factorization whose neighboring operators are
    positive semidefinite, and let $T_C$ be the trace matrix on its
    cyclic-active sectors. Assume that every ordered pair of cyclic-active
    sectors is joined by a path of nonzero neighboring operators. Assume also
    that, for every $L>0$, the normalized $L$-site marginals obtained from
    chain lengths $L+2$ and $L+1$ are equal, and that the closed operator has
    a positive real trace $Z_N$ at every length $N\geq2$. Then there is a
    number $\lambda>0$ such that
    \begin{align}
        (\lambda^{-1}T_C)^2
        &=(\lambda^{-1}T_C)^3.
        \notag
    \end{align}
\end{theorem}

\begin{proof}\leanok
    \uses{def:mpdo_suffix_sector_contraction,
      thm:mpdo_suffix_marginal_block_expansion,
      thm:mpdo_cyclic_active_square_positive_adjacent_marginals}
    By
    Theorem~\ref{thm:mpdo_cyclic_active_square_positive_adjacent_marginals},
    every entry of $T_C^2$ is positive. For fixed $q,h$, choose a positive
    summand $(T_C)_{h,r}(T_C)_{r,q}$ of $(T_C^2)_{h,q}$. The retained
    three-sector word $h,r,q$ has nonzero trace product. Comparing the two
    adjacent marginals at retained length three gives
    \begin{align}
        (T_C^3)_{q,h}
        &=\frac{Z_5}{Z_4}(T_C^2)_{q,h}.
        \notag
    \end{align}
    The ratio is independent of $q,h$. Taking $\lambda=Z_5/Z_4$ proves the
    assertion.
\end{proof}

% Source: arXiv:1606.00608, Appendix C.2, Proposition `4to2`, lines
% 1597--1619.
% **Scope restriction (coefficient extraction):** this proves the normalized
% square--cube identity preceding the source's rank-one conclusion. See
% docs/paper-gaps/cpsv16_commuting_form_to_sal.tex.
\begin{theorem}[Selected fixed-product square--cube identity]
    \label{thm:mpdo_selected_fixed_product_square_cube}
    \lean{MPOTensor.EtaLocalStructureData.%
      exists_normalized_selectedFixedProduct_cyclicActiveSectorTraceMatrix_pow_two_eq_pow_three}
    \leanok
    \uses{def:mpdo, def:injective, def:normal, def:mpo_source_zcl,
      def:mpdo_eta_local_structure_data,
      def:mpdo_physical_sector_factorization,
      thm:mpdo_eta_fixed_bond_product_tensor,
      thm:mpdo_fixed_bond_positive_physical_sector_representative,
      def:mpdo_cyclic_active_sector}
    Let $K$ be an injective normal matrix-product density tensor with source
    zero correlation length and an eta-local structure. Let $F$ be the
    positive physical-sector factorization of the selected fixed-product
    tensor, and let $T_C$ be its trace matrix on cyclic-active sectors. There
    is a number $\lambda>0$ such that
    \begin{align}
        (\lambda^{-1}T_C)^2=(\lambda^{-1}T_C)^3.
        \notag
    \end{align}
\end{theorem}

\begin{proof}\leanok
    \uses{lem:mpdo_source_zcl_marginal_stability,
      thm:mpdo_eta_fixed_bond_product_tensor,
      thm:mpdo_fixed_bond_positive_physical_sector_representative,
      thm:mpdo_selected_fixed_product_cyclic_recurrence,
      thm:mpdo_selected_fixed_product_adjacent_marginals,
      thm:mpdo_selected_fixed_product_positive_traces,
      thm:mpdo_cyclic_active_square_cube_adjacent_marginals}
    Express the original tensor in the physical coordinates of $F$. At every
    length $N\geq2$, its closed operator is a positive scalar multiple of the
    selected closed operator. Source zero correlation length therefore gives
    the adjacent marginal equality for the selected tensor, and positivity
    gives positive selected traces.
    Theorem~\ref{thm:mpdo_selected_fixed_product_cyclic_recurrence} supplies
    recurrence.
    Theorem~\ref{thm:mpdo_cyclic_active_square_cube_adjacent_marginals} now
    gives the required normalization.
\end{proof}

\begin{theorem}[Markov cuts from the selected fixed product]
    \label{thm:mpdo_selected_fixed_product_markov_cuts}
    \lean{MPOTensor.EtaLocalStructureData.%
      exists_hayashiMarkovDecomposition_selectedPhysicalCut_of_isSourceZCL}
    \leanok
    \uses{def:mpdo, def:injective, def:normal, def:mpo_source_zcl,
      def:mpdo_eta_local_structure_data,
      thm:mpdo_fixed_bond_positive_physical_sector_representative,
      def:hayashi_markov_decomposition}
    Let $K$ be an injective normal tensor generating matrix-product density
    operators, with an eta-local structure and source zero correlation
    length. Express $K$ in the one-site physical coordinates of the positive
    fixed-product factorization selected by the eta-local structure. For
    arbitrary $A,C\geq0$, the normalized marginal on consecutive regions of
    lengths $A$, $1$, and $C$, obtained from the periodic chain of length
    $A+C+2$, admits a Hayashi Markov decomposition.
\end{theorem}

\begin{proof}\leanok
    \uses{thm:mpdo_selected_fixed_product_cyclic_recurrence,
      thm:mpdo_selected_fixed_product_adjacent_marginals,
      thm:mpdo_selected_fixed_product_iterated_marginals,
      thm:mpdo_selected_fixed_product_positive_traces,
      thm:mpdo_selected_fixed_product_reduced_states,
      thm:mpdo_selected_fixed_product_square_cube,
      thm:mpdo_cyclic_active_square_positive_adjacent_marginals,
      thm:mpdo_cyclic_active_explicit_fourth_region_blocks,
      thm:mpdo_cyclic_active_markov_of_fourth_region,
      thm:matrix_primitive_pow_two_rank_one}
    The selected trace matrix is recurrent, and its square is strictly
    positive. After the positive normalization of
    Theorem~\ref{thm:mpdo_selected_fixed_product_square_cube}, its square and
    cube agree. Thus the normalized matrix is primitive, with its second
    power as a positive power, and the latter factors as
    $a b^{\mathsf T}$ for strictly positive vectors $a,b$.

    Transport the three-suffix/one-suffix marginal identity from the source
    tensor to the selected fixed product. In cyclic-active coordinates, its
    retained blocks are the unchanged interior neighboring products tensored
    with the boundary factors determined by $a$ and $b$. Regrouping at the
    distinguished middle site gives the direct sum of normalized positive
    left and right path states. Their trace products are nonnegative and sum
    to one. Finally, positive proportionality identifies this selected
    normalized marginal with the marginal of $K$ in the same physical
    coordinates.
\end{proof}

% **Local fix (cyclic-active restriction):** this is a restricted replacement for
% the primitive-matrix step of the SAL-dependent CPSV16 Lemma `propSN`. See
% docs/paper-gaps/cpsv16_commuting_form_to_sal.tex.
\begin{theorem}[Primitivity on cyclic-active sectors]
    \label{thm:mpdo_cyclic_active_trace_matrix_primitive}
    \lean{MPOTensor.PhysicalSectorFactorization.exists_sectorVirtualMatrix_ne_zero_of_isCyclicActiveSector}
    \lean{MPOTensor.PhysicalSectorFactorization.isCyclicActiveSector_of_sectorVirtualMatrix_ne_zero}
    \lean{MPOTensor.PhysicalSectorFactorization.sectorVirtualMatrix_eq_zero_of_not_isCyclicActiveSector}
    \lean{MPOTensor.PhysicalSectorFactorization.nonempty_cyclicActiveSector}
    \lean{MPOTensor.PhysicalSectorFactorization.cyclicActiveSectorOneSiteMatrixFamily_span_eq_top}
    \lean{MPOTensor.PhysicalSectorFactorization.cyclicActiveSector_exists_sectorVirtualMatrix_ne_zero}
    \lean{MPOTensor.PhysicalSectorFactorization.exists_cyclicActive_two_edge_return}
    \lean{MPOTensor.PhysicalSectorFactorization.cyclicActiveSectorTraceMatrix_isIrreducible}
    \lean{MPOTensor.PhysicalSectorFactorization.cyclicActiveSectorTraceMatrix_pow_two_diag_pos}
    \lean{MPOTensor.PhysicalSectorFactorization.cyclicActiveSectorTraceMatrix_pow_three_diag_pos}
    \lean{MPOTensor.PhysicalSectorFactorization.cyclicActiveSectorTraceMatrix_isPrimitive}
    \leanok
    \uses{def:injective, def:mpdo_cyclic_active_sector}
    Suppose that the virtual dimension is positive, the MPO tensor is
    injective, and every $\eta_{k,h}$ is positive semidefinite. Then $C$ is
    nonempty and the matrix $T$ is primitive. More precisely, $T$ is
    irreducible and every $k\in C$ satisfies
    \begin{align}
        (T^2)_{k,k}&>0,
        \notag\\
        (T^3)_{k,k}&>0.
        \notag
    \end{align}
\end{theorem}

\begin{proof}\leanok
    The virtual matrices of all physical sectors span the full matrix
    algebra. If $A=A_k^{xy}$ is nonzero, nondegeneracy of the trace pairing
    gives a virtual matrix $B=A_q^{uv}$ such that $\tr(AB)\ne0$.
    Cyclicity gives $\tr(BA)=\tr(AB)\ne0$, and therefore
    $\eta_{k,q}\ne0$ and $\eta_{q,k}\ne0$.
    Thus a sector containing a nonzero virtual matrix lies on a directed
    two-cycle. Hence every virtual matrix outside $C$ vanishes, and the
    virtual matrices belonging to $C$ still span the full algebra.

    If a proper directed cut separated two nonempty subsets $C_1,C_2$ of
    $C$, let $V_{C_i}$ denote the span of the virtual matrices in $C_i$.
    Absence of an edge across the cut gives
    \begin{align}
        V_{C_1}V_{C_2}&=\{0\},
        \notag\\
        V_{C_1}+V_{C_2}&=M_D(\mathbb C).
        \notag
    \end{align}
    Both spaces contain a nonzero matrix. This contradicts the one-sided
    product obstruction for two nonzero subspaces whose sum is the full
    matrix algebra. Thus $T$ is irreducible. The same trace pairing gives a
    two-edge return through each sector and closes every retained edge
    through a third retained sector. These yield the displayed positive
    diagonal entries of $T^2$ and $T^3$; the coprime return lengths imply
    primitivity.
\end{proof}

% **Local fix (cyclic-active restriction):** CPSV16 line 1613 uses the one-step
% coefficient, whereas this restricted theorem proves rank one for its square.
% See docs/paper-gaps/cpsv16_commuting_form_to_sal.tex.
\begin{theorem}[Rank-one cyclic-active two-step coefficient]
    \label{thm:mpdo_cyclic_active_zcl_rank_one}
    \lean{MPOTensor.PhysicalSectorFactorization.cyclicActiveLeftRestriction}
    \lean{MPOTensor.PhysicalSectorFactorization.cyclicActiveLeftRestriction_leftTensor}
    \lean{MPOTensor.PhysicalSectorFactorization.cyclicActiveLeftRestriction_rightTensor}
    \lean{MPOTensor.PhysicalSectorFactorization.cyclicActiveLeftRestriction_neighboringOperator}
    \lean{MPOTensor.PhysicalSectorFactorization.reindex_cyclicActiveLeftRestriction_activeSectorTraceMatrix}
    \lean{MPOTensor.PhysicalSectorFactorization.cyclicActiveSectorTraceMatrix_normalized_relations_of_isSourceZCL}
    \lean{MPOTensor.PhysicalSectorFactorization.cyclicActiveSectorTraceMatrix_rank_pow_two_eq_one}
    \lean{MPOTensor.PhysicalSectorFactorization.normalized_cyclicActiveSectorTraceMatrix_rank_pow_two_eq_one}
    \lean{MPOTensor.PhysicalSectorFactorization.exists_normalized_cyclicActiveSectorTraceMatrix_rank_pow_two_eq_one_of_isSourceZCL}
    \leanok
    \uses{def:mpo_source_zcl,
      thm:mpdo_cyclic_active_trace_matrix_primitive,
      thm:matrix_primitive_pow_two_rank_one}
    Under the hypotheses of
    Theorem~\ref{thm:mpdo_cyclic_active_trace_matrix_primitive}, suppose that
    the physical-trace transfer has source zero correlation length. There is
    a number $\lambda>0$ such that, with $\widehat T=\lambda^{-1}T$,
    \begin{align}
        \widehat T^2&=\widehat T^3,
        \notag\\
        \rank(\widehat T^2)&=1,
        \notag
    \end{align}
    and $\tr(\widehat T^m)=\tr(\widehat T)$ for every $m\geq1$.
    By matrix multiplication, the two-step coefficient is
    $(\widehat T^2)_{k,h}=\sum_{q\in C}\widehat T_{k,q}\widehat T_{q,h}$. The assertion neither identifies
    $\widehat T$ with $\widehat T^2$ nor implies that the unreduced Beigi trace
    matrix is primitive or rank one.
\end{theorem}

\begin{proof}\leanok
    Set the left sector tensor equal to zero outside $C$. The preceding span
    argument shows that this does not change the physical tensor, and the
    neighboring operators with both indices in $C$ are unchanged. The
    source zero-correlation-length relation, transported through this
    restricted factorization, gives
    $\widehat T^2=\widehat T^3$ on the cyclic-active indices. By
    Theorem~\ref{thm:mpdo_cyclic_active_trace_matrix_primitive},
    $\widehat T$ is primitive. Its square is therefore a strictly positive
    idempotent matrix and has rank one.
\end{proof}

% Tracked in https://github.com/LionSR/TNLean/issues/4445.
% **Local fix (cyclic-active restriction):** the intermediate sum is over the
% positive-length cyclic support.  See
% docs/paper-gaps/cpsv16_commuting_form_to_sal.tex.
\begin{theorem}[Cyclic-active three-boundary trace coefficient]
    \label{thm:mpdo_cyclic_active_three_boundary_trace}
    \lean{MPOTensor.PhysicalSectorFactorization.%
      sum_cyclicActive_trace_mul_trace_eq_pow_two}
    \lean{MPOTensor.PhysicalSectorFactorization.%
      neighboringOperator_trace_eq_cyclicActiveSectorTraceMatrix}
    \lean{MPOTensor.PhysicalSectorFactorization.%
      sum_cyclicActive_normalized_trace_mul_trace_eq_pow_two}
    \lean{MPOTensor.PhysicalSectorFactorization.%
      sum_cyclicActive_partialTraceRight_threeSiteSectorClosure}
    \leanok
    \uses{def:mpdo_cyclic_active_sector,
      thm:mpdo_sector_closure_partial_traces}
    Let $C$ be the set of cyclic-active sectors and suppose that every
    neighboring operator is positive semidefinite. For $q,h\in C$,
    \begin{align}
        \sum_{r\in C}\tr(\eta_{q,r})\tr(\eta_{r,h})
        &=(T_C^2)_{q,h}.
        \notag
    \end{align}
    More generally, if $B_{q,h}(X)$ denotes the outer boundary operator in
    the three-site closure, then
    \begin{align}
        \sum_{r\in C}\tr_{\mathrm{mid}}
        (B_{q,h}(X)\otimes\eta_{q,r}\otimes\eta_{r,h})
        &=(T_C^2)_{q,h}B_{q,h}(X).
        \notag
    \end{align}
    Here $\tr_{\mathrm{mid}}$ traces both neighboring factors. No equality
    between $T_C$ and $T_C^2$ is asserted.

    For every $\lambda\in\R$, set
    $\widehat T_C=\lambda^{-1}T_C$. The same calculation with each fully
    traced edge multiplied by $\lambda^{-1}$ gives
    \begin{align}
        \sum_{r\in C}
        \lambda^{-1}\tr(\eta_{q,r})
        \lambda^{-1}\tr(\eta_{r,h})
        &=(\widehat T_C^2)_{q,h}.
        \notag
    \end{align}
\end{theorem}

\begin{proof}\leanok
    Positivity makes each $\tr(\eta_{a,b})$ real. Expanding matrix
    multiplication gives
    \begin{align}
        (T_C^2)_{q,h}
        &=\sum_{r\in C}(T_C)_{q,r}(T_C)_{r,h}
         =\sum_{r\in C}\tr(\eta_{q,r})\tr(\eta_{r,h}).
        \notag
    \end{align}
    The partial trace of the fixed-sector three-site closure is
    $\tr(\eta_{q,r})\tr(\eta_{r,h})B_{q,h}(X)$.
    Summing over $r\in C$ proves the second identity. Multiplying each of the
    two trace factors by $\lambda^{-1}$ gives the normalized identity.
\end{proof}

\begin{theorem}[Fourth-region blocks from an explicit factorization]
    \label{thm:mpdo_cyclic_active_explicit_fourth_region_blocks}
    \lean{MPOTensor.PhysicalSectorFactorization.%
      reindex_reducedBlockState_add_three_eq_cyclicActiveFourthRegionBlock}
    \leanok
    \uses{def:mpdo_physical_sector_factorization,
      def:mpdo_cyclic_active_sector}
    Let $F$ be a physical-sector factorization with positive semidefinite
    neighboring operators. Suppose that $\lambda>0$ and that functions
    $a,b$ on the cyclic-active sectors satisfy
    \[
        (\lambda^{-1}T_C)^2=a b^{\mathsf T}.
    \]
    For every $n\geq0$, reindex the $(n+1)$-site reduced state obtained from
    the periodic chain of length $n+4$ by its retained open-edge coordinates.
    The result is the inverse trace of the length-$(n+4)$ periodic operator
    times the direct sum over retained sector words. A word outside cyclic
    support contributes zero; every cyclic-active word contributes its
    unchanged interior neighboring product tensored with the separated
    boundary factors determined by $\lambda,a,b$.
\end{theorem}

\begin{proof}\leanok
    \uses{thm:mpdo_cyclic_active_three_boundary_trace,
      thm:mpdo_cyclic_active_deletion}
    Expand the reduced state as a direct sum of three-suffix contractions.
    Terms outside positive-length cyclic support vanish. For a retained
    cyclic-active word, the three traced suffix sites leave the interior
    product and the two-step boundary contraction. Substitute
    $(\lambda^{-1}T_C)^2=a b^{\mathsf T}$ to separate its two boundary
    factors, retaining the normalization by the length-$(n+4)$ trace.
\end{proof}

\begin{theorem}[Cyclic-active fourth-region factorization]
    \label{thm:mpdo_cyclic_active_fourth_region_factorization}
    \lean{MPOTensor.PhysicalSectorFactorization.%
      exists_cyclicActiveFourthRegion_formula_of_isSourceZCL}
    \leanok
    \uses{thm:mpdo_source_zcl_iterated_marginal_stability,
      thm:mpdo_cyclic_active_explicit_fourth_region_blocks,
      thm:mpdo_cyclic_active_three_boundary_trace,
      thm:mpdo_cyclic_active_deletion}
    Let $\mathcal K$ be injective and have source zero correlation length,
    and suppose its neighboring operators are positive semidefinite. There
    are $\lambda>0$ and strictly positive functions $a,b$ on the
    cyclic-active sectors, normalized by $a\mathbin{\boldsymbol\cdot}b=1$,
    such that every fourth-region marginal is a direct sum over retained
    sector words of the unchanged interior neighboring product tensored with
    the separated left and right boundary factors determined by $a$ and $b$.

    % Local fix (cyclic-active restriction): the coefficient being factorized
    % is $(\lambda^{-1}T_C)^2$, obtained from one additional source-ZCL
    % marginal replacement and three traced boundary sites. No equality
    % between $T_C$ and $T_C^2$ is used. This correction is recorded in
    % docs/paper-gaps/cpsv16_commuting_form_to_sal.tex.
\end{theorem}

\begin{proof}\leanok
    The additional marginal replacement leaves the retained interior product
    unchanged and produces the two-step coefficient on the cyclic-active
    sectors. Its positive rank-one factorization separates the two surviving
    boundary sums. All sector words leaving the cyclic-active support vanish
    before the two outer partial traces are evaluated.
\end{proof}

\begin{theorem}[Markov decomposition from an explicit fourth-region formula]
    \label{thm:mpdo_cyclic_active_markov_of_fourth_region}
    \lean{MPOTensor.PhysicalSectorFactorization.%
      exists_hayashiMarkovDecomposition_cyclicActiveCut_of_fourthRegion}
    \leanok
    \uses{def:mpdo_physical_sector_factorization,
      def:mpdo_cyclic_active_sector,
      def:hayashi_markov_decomposition}
    Let $F$ be a positive-dimensional physical-sector factorization whose
    neighboring operators are positive semidefinite. Fix $A,C\geq0$,
    $\lambda\in\mathbb R$, and strictly positive functions $a,b$ on the
    cyclic-active sectors. Suppose that the length-$(A+C+4)$ periodic
    operator has positive trace, the length-$(A+C+2)$ periodic operator has
    nonzero trace, and the normalized $(A+C+1)$-site reduced state at length
    $A+C+2$ has the explicit fourth-region block formula determined by
    $\lambda,a,b$. Then its tripartite form on regions of lengths
    $A$, $1$, and $C$ admits a Hayashi Markov decomposition.
\end{theorem}

\begin{proof}\leanok
    Regroup every retained fourth-region block at the distinguished middle
    sector. Its positive left and right path factors are normalized by their
    traces, with fixed density matrices used only when a factor vanishes.
    The positive long-chain trace makes all block weights nonnegative. The
    nonzero short-chain trace makes the reduced state trace one, so these
    weights sum to one. The resulting direct sum is exactly the Hayashi
    decomposition in the sector-coordinate middle-site splitting.
\end{proof}

\begin{theorem}[Cyclic-active Markov decomposition at every cut]
    \label{thm:mpdo_cyclic_active_markov_all_cuts}
    \lean{MPOTensor.PhysicalSectorFactorization.%
      exists_hayashiMarkovDecomposition_cyclicActiveCut_of_isSourceZCL}
    \leanok
    \uses{thm:mpdo_cyclic_active_fourth_region_factorization,
      thm:mpdo_cyclic_active_markov_of_fourth_region,
      def:hayashi_markov_decomposition}
    Under the hypotheses of
    Theorem~\ref{thm:mpdo_cyclic_active_fourth_region_factorization}, for
    arbitrary $A,C\geq 0$, the fourth-region marginal on consecutive regions
    of lengths $A$, $1$, and $C$ admits a Hayashi Markov decomposition.
    The middle site decomposes as the direct sum of its left and right sector
    factors. Each block is the tensor product of normalized positive left
    and right path states, and its probability is the product of their
    unnormalized traces divided by the positive periodic normalization.

    % Local fix (cyclic-active restriction): the block weights and path states
    % arise from the factorization of $(\lambda^{-1}T_C)^2$ described above.
    % This correction is recorded in
    % docs/paper-gaps/cpsv16_commuting_form_to_sal.tex.
\end{theorem}

\begin{proof}\leanok
    Separate the retained chain at the distinguished middle site. The
    retained neighboring factors split into the left and right path products,
    while the two cyclic endpoint factors become their respective boundary
    matrices. Normalize each positive block by its trace, using an arbitrary
    trace-one positive matrix when the block vanishes. The trace-one identity
    for the full marginal shows that the resulting non-negative block weights
    sum to one.
\end{proof}

\begin{theorem}[Physical-sector factorization and ZCL imply SAL]
    \label{thm:mpdo_physical_sector_factorization_zcl_sal}
    \lean{MPOTensor.PhysicalSectorFactorization.isSAL_of_isSourceZCL}
    \leanok
    \uses{thm:mpdo_cyclic_active_markov_all_cuts,
      thm:mpdo_all_cut_markov_implies_sal,
      thm:mpdo_physical_support_restriction_sal}
    Let $\mathcal K$ have positive bond dimension and be injective, and let
    $F$ be a physical-sector factorization of $\mathcal K$. Suppose every
    neighboring operator $\eta^F_{q,h}$ is positive semidefinite. If the
    physical-trace transfer of $\mathcal K$ has source zero correlation
    length, then $\mathcal K$ saturates the area law.

    % **Scope restriction (physical-sector factorization):** The theorem
    % assumes the factorization F and positivity of every neighboring
    % operator.  CPSV16 Proposition 4to2, lines 1597--1619, starts instead
    % from a translation-invariant commuting bond and the original source
    % tensor.  The selected fixed tensor has the required positive
    % factorization, and
    % Theorem~\ref{thm:mpdo_selected_fixed_product_cyclic_recurrence}
    % establishes recurrence through the auxiliary space of the original
    % tensor.  It remains to derive the rank-one Markov conclusion for the
    % original source without imposing source zero correlation length on the
    % selected tensor; see docs/paper-gaps/cpsv16_commuting_form_to_sal.tex.
\end{theorem}

\begin{proof}\leanok
    Positivity of the neighboring operators makes every nonempty
    sector-coordinate chain operator positive. Source zero correlation
    length makes its trace nonzero. Apply the preceding theorem with
    $A=m-1$ and $C=N-m-1$ at every admissible cut, then use the all-cut
    Markov criterion. Finally transport saturation through the physical
    isometry of $F$.
\end{proof}

\begin{theorem}[A translation-invariant commuting bond and ZCL imply SAL]
    \label{thm:mpdo_translation_invariant_commuting_form_zcl_sal}
    \notready
    \uses{def:mpdo_eta_local_structure_data, def:mpo_source_zcl,
        def:is_sal, def:mpdo, def:mpdo_injective_inverse_layer}
    Let $\mathcal K$ be an injective normal tensor generating MPDOs. Suppose
    that there is one
    positive two-site bond $B$ whose translates commute and, for every
    $N\geq2$, a constant $c_N>0$ such that
    \begin{align}
        \sigma^{(N)}(\mathcal K)
        &=c_N\prod_{i=0}^{N-1}B_{i,i+1}.
        \notag
    \end{align}
    If the physical-trace transfer of $\mathcal K$ has zero correlation
    length, then $\mathcal K$ saturates the area law.

    This is the proposition at
    \cite[Appendix~C.2, lines~1597--1619]{Cirac2016MPDO_arXiv}.
    % The local positive neighboring operators are constructed.  The printed
    % rank-one inference at source line 1613 is obstructed by Lemma C.5. Gap
    % documented in
    % docs/paper-gaps/cpsv16_commuting_form_to_sal.tex.
\end{theorem}

\begin{theorem}[$\eta$-local structure gives the one-sector GSNNCH form]
    \label{thm:mpdo_eta_local_structure_has_gsnnch_form}
    \lean{MPOTensor.EtaLocalStructureData.hasGSNNCHForm}
    \lean{MPOTensor.hasGSNNCHForm_of_etaLocalStructure}
    \leanok
    \uses{thm:mpdo_eta_local_structure_has_commuting_bond_product,
        thm:mpdo_commuting_bond_product_has_gsnnch_form}
    The single positive bond determined by the $\eta$-local structure gives
    the one-sector instance of the explicit GSNNCH decomposition.
\end{theorem}

\begin{proof}\leanok
    Apply the one-sector construction to the single-bond product at every
    chain length.
\end{proof}

\begin{theorem}[Injective SAL tensors generate GSNNCH states]
    \label{thm:mpdo_injective_sal_is_gsnnch}
    \lean{MPOTensor.isGSNNCH_of_isInjective_of_isSAL}
    \leanok
    \uses{def:is_sal, cor:mpdo_eta_local_structure_of_is_sal,
        thm:mpdo_eta_local_structure_has_gsnnch_form,
        thm:mpdo_gsnnch_of_sector_form}
    Every injective MPO tensor satisfying SAL generates Gibbs states of a
    nearest-neighbor commuting Hamiltonian: the normalized finite-chain
    states are density operators with the explicit sector form of
    Definition~\ref{def:mpdo_source_gsnnch}, realized through one sector.
    The one-sector form is the commuting product form in which
    \cite[Proposition~C.8]{Cirac2016MPDO_arXiv} states its conclusion for an
    injective tensor, combined here with the density normalization of
    \cite[Definition~4.8]{Cirac2016MPDO_arXiv}.
\end{theorem}

\begin{proof}\leanok
    Corollary~\ref{cor:mpdo_eta_local_structure_of_is_sal} gives the
    $\eta$-local structure, whose single positive bond gives the one-sector
    form by Theorem~\ref{thm:mpdo_eta_local_structure_has_gsnnch_form}. SAL
    includes the nonvanishing of every positive-length trace, so
    Theorem~\ref{thm:mpdo_gsnnch_of_sector_form} applies.
\end{proof}

\begin{remark}[Positive physical-sector factorization]
    Suppose that $K$ satisfies $\mathrm{SAL}(K)$ and admits a physical-sector
    factorization whose neighboring operators $\eta_{k,h}$ are simultaneously
    positive semidefinite. The associated two-site bond is
    \begin{align}
        B&=\sum_{k,h}
        \Id_{H_L^{(k)}}\otimes
        \eta_{k,h}\otimes
        \Id_{H_R^{(h)}},
        \label{eq:mpdo_eta_bond_assembly}\\
        B&\geq0,
        \notag
    \end{align}
    and, for every $N\geq2$,
    \begin{align}
        \sigma^{(N)}(K)
        &=c_N\prod_{n=1}^{N}B_{n,n+1},
        \label{eq:mpdo_eta_product_form}\\
        c_N&>0,
        \notag\\
        [B_{i,i+1},B_{j,j+1}]&=0.
        \notag
    \end{align}
    Under this hypothesis,
    Theorem~\ref{thm:mpdo_eta_local_structure_of_positive_physical_sector_factorization}
    gives the $\eta$-local structure and proves
    (\ref{eq:mpdo_eta_bond_assembly})--(\ref{eq:mpdo_eta_product_form}) with
    $c_N=1$. Theorem~\ref{thm:mpdo_eta_local_structure_has_commuting_form}
    then gives the commuting-form conclusion. Adjacent translates commute on every
    periodic chain of length at least three by
    Theorem~\ref{thm:mpdo_periodic_adjacent_physical_bonds_comm}, and the
    crossed two-site case is
    Theorem~\ref{thm:mpdo_physical_two_site_crossed_bonds_comm}. Disjoint
    translates commute by
    Theorem~\ref{thm:mpdo_periodic_disjoint_physical_bonds_comm}. These cases
    give pairwise commutativity at every length $N\geq2$ by
    Theorem~\ref{thm:mpdo_periodic_pairwise_physical_bonds_comm}. Given the
    physical-sector factorization, the all-length product identity is
    Theorem~\ref{thm:mpdo_exact_physical_sector_bond_product}.
    The doubled-index transfer condition enters only in the subsequent RFP and
    single-bond branch conclusions.
\end{remark}
